% ECONOMICS_PROBLEM: {"id": "R23-505", "title": "Gentrification and incumbent welfare", "jel_code": "R23", "source_id": "OP-0465"}
% FIRST_POSTED: 2026-10-09
% ATTEMPT_STATUS: partial
% ENTRY_KIND: research_attempt
\documentclass[11pt,a4paper]{article}
\usepackage[T1]{fontenc}
\usepackage[utf8]{inputenc}
\usepackage[margin=27mm]{geometry}
\usepackage{amsmath,amssymb,amsthm,mathtools}
\usepackage[hidelinks]{hyperref}
\setlength{\parindent}{0pt}
\setlength{\parskip}{5pt}
\newtheorem{theorem}{Theorem}[section]
\newtheorem{proposition}[theorem]{Proposition}
\newtheorem{lemma}[theorem]{Lemma}
\newcommand{\R}{\mathbb R}
\newcommand{\E}{\mathbb E}
\DeclareMathOperator{\Var}{Var}
\DeclareMathOperator{\Cov}{Cov}
\begin{document}
\section*{Incumbent welfare under revitalization: choices, ties, and displacement}
\section{Target and partial result}
A neighborhood's mean income can rise because its original residents gain, because high-income households replace them, or both. The catalogue instead follows households present before announcement. This attempt gives a sharp finite-choice welfare calculation with uncertain social ties, a transparent displacement-classification sensitivity formula, and bounds for missing movers. None supplies actual ten-year treatment effects without the program and cohort data.

Normalize all amounts to discounted baseline dollars over the stated ten-year horizon. In a restricted quasi-linear model, a baseline renter chooses between staying, $h$, and moving, $o$. Let $v_a$ be consumption-equivalent benefit from staying under regime $a$, already net of rent, travel and other measured costs. Let $b_a$ be the best outside option net of moving costs. A persistent individual tie value $t\in[0,T]$ is received only at $h$. The indirect money value is
\[
 V_a(t)=\max\{v_a+t,b_a\},\qquad
 \Delta(t)=V_1(t)-V_0(t). \tag{1}
\]
The same $t$ enters both regimes; bounding each maximum independently would discard a real cross-regime restriction. Quasi-linearity, fixed outside-option summaries and transferable money are modeling assumptions. They do not assert that every social relationship has a measurable market price.

\section{Exact sharp welfare bounds with a common tie value}
Define thresholds $t_a=b_a-v_a$ and the finite set
\[
 \mathcal K=\{0,T\}\cup\{t_0,t_1:\ 0\le t_a\le T\}.
\]
\begin{proposition}
The exact identified interval for the individual's welfare change under (1) is
\[
 [\min_{t\in\mathcal K}\Delta(t),\quad\max_{t\in\mathcal K}\Delta(t)]. \tag{2}
\]
If observed optimizing choices impose an additional nonempty closed interval $[L,U]\subseteq[0,T]$, replace $0,T$ by $L,U$ and retain only thresholds in that interval.
\end{proposition}
\begin{proof}
Each maximum in (1) is continuous and piecewise affine, with its only kink at $t_a$. Their difference is affine between consecutive points in $\mathcal K$. Every extremum on one such interval occurs at an endpoint, unless the function is constant there. This proves (2). Continuity and connectedness of $[0,T]$ imply that every intermediate value is attained. Every admissible $t$ is a compatible primitive, establishing sharpness. The same reasoning applies after intersecting the support with choice inequalities.
\end{proof}
Under weak optimality, observing stay in regime $a$ imposes $t\ge t_a$ and observing move imposes $t\le t_a$. A specified tie-breaking rule may make a boundary open; in that case (2) reports sharp infimum and supremum, with endpoint attainment recorded separately. Observing both potential choices for the same household is normally impossible, so an empirical restriction cannot condition on an unobserved choice.

For an arithmetic example let $v_0=0,b_0=2,v_1=-1,b_1=4,T=6$. Before the program the household stays when $t\ge2$; afterward it stays when $t\ge5$. Welfare changes equal $2$ for $t\le2$, $4-t$ for $2\le t\le5$, and $-1$ for $t\ge5$. Thus the sharp welfare interval is $[-1,2]$. A household with $t=3$ moves and gains one; one with $t=4.5$ moves and loses one-half. Moving alone does not determine the sign of welfare, and a rent increase can harm a high-tie incumbent even when a better outside option becomes available.

\section{More locations: a finite linear-program method}
For $J$ possible destinations, let
$V_a(\theta)=\max_j\{c_{aj}+d_{aj}'\theta\}$, where uncertain preference vector $\theta$ belongs to a nonempty compact polytope $\Theta$. Enumerate pairs $(j,k)$ of optimizing destinations under regimes one and zero. In each pair's region impose
\[
 c_{1j}+d_{1j}'\theta\ge c_{1l}+d_{1l}'\theta\quad\forall l,
 \qquad
 c_{0k}+d_{0k}'\theta\ge c_{0l}+d_{0l}'\theta\quad\forall l.
\]
On this region the welfare difference is the affine function
$c_{1j}-c_{0k}+(d_{1j}-d_{0k})'\theta$. Minimize and maximize it by linear programming, skipping empty regions, and take the extrema across all $J^2$ pairs. Every $\theta$ has at least one optimizing pair, including at ties, so this union covers the complete problem. Conversely every feasible pair-region point evaluates the original maxima exactly. Compactness guarantees attainable extrema. The procedure is therefore exact for this finite-affine benchmark, not a numerical guess about an arbitrary household utility function.

For a finite baseline population with separately bounded independent household primitives and fixed weights, the sharp average interval is the weighted sum of household intervals. Shared parameters, budget restrictions, or equilibrium rents couple households and can make that summation nonsharp. They must instead be included in one joint optimization. Landlords' gains should be reported separately or included through the declared ownership weights; counting both rent receipts and the capitalization of the same receipts as two contemporaneous benefits is incorrect.

\section{Operational displacement and classification uncertainty}
Let $D$ be the predeclared operational event, such as an eviction or a move following an excessive rent burden, and let $\widetilde D$ be a fallible recorded flag. Under constant sensitivity $s=\Pr(\widetilde D=1\mid D=1)$ and specificity $c=\Pr(\widetilde D=0\mid D=0)$ within a regime,
\[
 o:=\Pr(\widetilde D=1)=s p+(1-c)(1-p),\quad
 p=\frac{o+c-1}{s+c-1}, \tag{3}
\]
provided $s+c>1$. This follows from total probability. Feasible $(s,c)$ additionally require the right side to lie in $[0,1]$. With known validation intervals for $s,c$, minimizing and maximizing the displayed continuous ratio over the feasible set gives sharp probability bounds, provided the denominator is bounded away from zero. Sharpness follows by constructing the corresponding two-by-two confusion table. Without a positive denominator bound, use the original linear probability equation and report possibly uninformative limits rather than divide by a nearly zero number.

For two regimes, opposite-endpoint subtraction is sharp if their classification parameters can vary independently. A common validation parameter couples the two probabilities and calls for joint optimization. Even perfect classification of an operational indicator does not reveal a person's unobserved causal reason for moving; the average policy effect on the defined indicator remains the estimand.

\section{Follow-up and the remaining empirical work}
Suppose a credible market-level assignment design identifies the regime-specific distribution among followed baseline households, with follow-up rate $r_a$, mean $m_a$ and a justified bound $V_a\in[\ell,u]$. Then
$E[V_a]\in[r_am_a+(1-r_a)\ell,r_am_a+(1-r_a)u]$. Assigning missing values to endpoints proves sharpness without assumptions on why households disappear. A household relocating outside the initial neighborhood is not missing if the cohort system follows it. New residents do not replace missing incumbents in this formula.

A full application still needs the anticipation date, the investment package, appropriate control neighborhoods, destination outcomes, and a defensible monetary-equivalent convention. Adjacent-area spillovers can violate a comparison that treats neighboring places as unaffected controls. The finite-choice calculations show exactly what a chosen preference model would imply and what tie uncertainty leaves open; they do not settle the catalogue's actual incumbent welfare or displacement effects.
\begin{thebibliography}{9}
\bibitem{ck} E. Chyn and L. F. Katz (2021), \emph{Neighborhoods Matter: Assessing the Evidence for Place Effects}, Journal of Economic Perspectives 35(4), 197--222. \url{https://www.aeaweb.org/articles?id=10.1257/jep.35.4.197}. Source motivation for incumbent and place-effect questions, rather than numerical inputs to this example.
\end{thebibliography}

\end{document}
